\documentclass[aspectratio=169]{beamer}

\usetheme{Madrid}
\usecolortheme{default}

\usepackage{booktabs}
\usepackage{hyperref}
\usepackage{xcolor}
\usepackage{array}
\usepackage{graphicx}

% Global bullet spacing
\newlength{\myitemsep}
\setlength{\myitemsep}{1em}
\let\olditemize\itemize
\renewcommand{\itemize}{\olditemize\setlength{\itemsep}{\myitemsep}}
\let\oldenumerate\enumerate
\renewcommand{\enumerate}{\oldenumerate\setlength{\itemsep}{\myitemsep}}

% ============================================================
%  HOUSE STYLE: see AI-Pulse/SLIDE_STYLE.md
%    No em-dashes (use colons). No boldface. No closing punchlines.
%    4-5 bullets a slide, one line each where possible. American spelling.
%    Citations on the References slide; source comments on every claim.
%
%  SLIDE STATUS LEGEND
%    (untagged)  settled, reviewed with Bruno
%    % [NEW]     built, not yet walked through
%    % [STAGE]   a stage of the running demo
%    % TODO:     not built yet
%
%  RUNNING DEMO: "How do I keep up with AI?" threaded through the session.
%    Shared capabilities are demoed once. Divergent ones are demoed twice.
%    Opening joke, slide 4: "A lot of people are worried about losing their jobs.
%    Today I will show you how to replace me entirely." Closed on slide 44.
% ============================================================

\title{The Major AI Ecosystems}
\subtitle{D-Lab AI Pulse, Fall 2026: Session 2}
\author[D-Lab]{Materials: \url{https://github.com/dlab-berkeley/AI_Pulse}\\Workshop suggestions: \url{https://forms.gle/x8KfejgCViE7t6iQA}}
\date{15 September 2026}

\begin{document}

% [NEW]
\begin{frame}
    \titlepage
\end{frame}

% [NEW]
\begin{frame}{Where Session 1 Left Off}
    \begin{itemize}
        \item AI has developed well past the chat box: agents, projects, and connectors
        \item Getting the most out of any of it takes an application, and a subscription
        \item Which raises three questions:
            \vspace{1em}
            \begin{itemize}\setlength{\itemsep}{4pt}
                \item What is the difference between the environments
                \item What is the difference between the plans
                \item Is it worth paying more, given what I actually do
            \end{itemize}
        \item Today is about answering all three
    \end{itemize}
\end{frame}

% [NEW]
% [NEW]
% OPENING JOKE goes here, delivered not printed:
%   "A lot of people are worried about losing their jobs.
%    Today I will show you how to replace me entirely."
% The running task is genuinely the job this workshop does for the room every two weeks,
% which is what makes the joke work. Loop closed on the "Did It Replace Me?" slide.
% The running task is now the last bullet rather than the title, 14 Sept.
\begin{frame}{Today's Roadmap}
    \begin{itemize}\setlength{\itemsep}{7pt}
        \item We look at the two main environments: OpenAI and Anthropic
        \item Part 1, the basics, what each company sells and what it costs
        \item Part 2, what they have in common, which is most of it
        \item Part 3, where they differ in practice, run side by side
        \item Part 4, what you are agreeing to, and what happens when it changes
        \item Throughout, one task carried end to end: how do I keep up with AI
    \end{itemize}
\end{frame}


% ============================================================
%  PART 1
% ============================================================

% [NEW]
\begin{frame}[plain]
    \vfill
    \begin{center}
        {\Huge Part 1}\\[0.8em]
        {\Large The Basics}
    \end{center}
    \vfill
\end{frame}

% [NEW]
% Bruno's prepared answer, agreed 1 Sept 2026. See google_question.md.
% Do NOT say "Google has fallen behind": Session 1 slide 13 has Gemini 3.7 Flash ninth on
% LMArena and slide 8 leads with Gemini's billion users. The deck would contradict itself.
% [NEW]
% Rewritten 7 Sept on Bruno's note. The old version argued from Berkeley's Gemini
% licence alone. This argues from what makes an environment worth committing to.
% Alternatives named: xAI's Grok, Moonshot's Kimi, Meta, Mistral, DeepSeek, Z.ai.
% !! Kimi K3 is open weights, so "general use" is doing different work there: you can
%    run it yourself, but there is no consumer subscription of the kind this session is
%    comparing. Say that if anyone asks.
\begin{frame}{Why These Two}
    \begin{itemize}
        \item Berkeley gives you Gemini and Gemini Notebook already, so that decision is
              made for you
        \item They are close to the frontier at all points
        \item Both build for general use, not for one narrow job, so what you learn
              transfers
        \item Both have a large community, which is where the workflows actually come from
        \item And both are established enough to be, somewhat, trustworthy with your work
        \item There are others for general use: Grok, Kimi, Mistral, DeepSeek
    \end{itemize}
\end{frame}


% [NEW]
% UPDATED 6 Sept 2026: Fable 5.1 and Mythos 5.1 shipped 1 Sept 2026. Fable 5.1 is generally
%   available at $10/$50 with cache reads cut 75% to $0.25. Mythos 5.1 is US organizations
%   only, via the verification programs. GPT-6 Astra launched 3 Sept but is a STAGED release
%   (alpha testers, then Daybreak Blue), so it is not something this room can buy: leave it
%   off the slide. See Ideas/Fall-2026/02_AI_Ecosystems/mythos_lsvp.md.
% TABLE REBUILT 14 Sept. !! CHECK EVERY NAME AND ROW BEFORE DELIVERY.
% Anthropic: Fable 5.1 and Mythos 5.1 (1 Sept 2026, $10/$50, cache reads $0.25).
%   Mythos 5.1 is US organizations only, via the Cyber and Life Sciences Verification
%   Programs. Opus 5 (24 Jul, $5/$25). Sonnet 5 (30 Jun, $2/$10 rising to $3/$15 on
%   1 Sept 2026). Haiku 4.5 (1 Oct 2025, 200k, $1/$5). No Haiku 5 exists.
% OpenAI: GPT-6 Astra shipped 3 Sept 2026 and rolled out 3-5 Sept to Pro, Enterprise and
%   Business Premium, then Plus and Business, plus the API and AWS. So it IS buyable now.
%   GPT-5.6 Sol / Terra / Luna, all 9 Jul 2026, 1.05M context.
% The o-series is gone (o4-mini, GPT-4o, 4.1, 5, 5.1, 5.2, 4.5 retired Feb-Jun 2026; o3
%   retired 26 Aug) but the slide no longer says so: the table carries the point.
% !! The tier labels (Frontier / Default / Balanced / Cheap) are OUR framing, inferred
%    from price. Neither vendor publishes them. Check they are defensible or relabel.
\begin{frame}{Two Model Lineups}
    \begin{itemize}
        \item Anthropic keeps parallel tiers. OpenAI ships variants of one model
    \end{itemize}

    \vspace{0.8em}
    \begin{center}
    \small
    \begin{tabular}{lll}
        \toprule
        & Anthropic & OpenAI \\
        \midrule
        Frontier      & Fable 5.1   & GPT-6 Astra \\
        Default       & Opus 5      & GPT-5.6 Sol \\
        Balanced      & Sonnet 5    & GPT-5.6 Terra \\
        Cheap         & Haiku 4.5   & GPT-5.6 Luna \\
        \bottomrule
    \end{tabular}
    \end{center}

    \vspace{0.8em}
    \begin{itemize}
        \item How hard the model thinks is now a slider, not a separate family of models
    \end{itemize}
\end{frame}



% [NEW] 14 Sept, on Bruno's note: the price table says what it costs, not what you get.
% Sources: Ideas/Fall-2026/02_AI_Ecosystems/ecosystems_reference.md sections 3 and 4.
% Anthropic Free is chat only: NO Claude Code, NO Cowork. Projects: Free 5, paid unlimited,
%   RAG mode on paid expands effective capacity ~10x.
% OpenAI Free has Codex and ChatGPT Work on desktop. Projects unlimited on all plans; files
%   capped Free 5 / Go-Plus 25 / Pro-Business-Edu 40.
% Skills: Anthropic on EVERY plan including Free, open SKILL.md standard. OpenAI Skills are
%   Business and Enterprise only; Custom GPTs need a paid plan.
% Scheduled jobs: Anthropic Routines, research preview, per-account daily cap, no count
%   limit published. OpenAI Tasks hard count limits: Go 3 / Plus 5 / Business 10 / Pro 15.
% !! CONTESTED, reference doc 12.2: "ChatGPT Work on Free" and "unlimited Luna text on
%    Free". The OpenAI Free cell depends on the first. Verify in a logged-out browser
%    before delivery, or soften that cell.
% Team rows dropped 14 Sept: both are $25/seat and the row added little.
% !! The 5x and 20x figures are the vendors' own advertised multipliers. BOTH ARE
%    CONTESTED: see the very next slide, and Ideas/.../plan_limits_dispute.md. This table
%    states what is advertised; the next slide says what it means.
\begin{frame}{What You Get, Not What You Pay}
    \begin{center}
    \small
    \begin{tabular}{p{1.5cm}>{\raggedright\arraybackslash}p{5.4cm}>{\raggedright\arraybackslash}p{5.0cm}}
        \toprule
        & Anthropic & OpenAI \\
        \midrule
        Free    & Chat, skills, 5 projects              & Chat, Codex, Work, projects \\
        \$8     & none                                  & + more files, 3 scheduled jobs \\
        \$20    & + Claude Code, Cowork, projects  & + 5 scheduled jobs \\
        \$100   & + five times the \$20 limits          & + five times, 15 scheduled jobs \\
        \$200   & + twenty times the \$20 limits        & + twenty times \\
        \bottomrule
    \end{tabular}
    \end{center}

    \vspace{0.8em}
    \begin{itemize}
        \item OpenAI has one rung Anthropic does not: an \$8 tier, which may carry ads
    \end{itemize}
\end{frame}

% [NEW]
% Anthropic: rolling 5-hour session + weekly caps (one all-model, one Sonnet-only).
%   Usage across claude.ai, Claude Code and Desktop draws from ONE SHARED POOL.
% OpenAI: four vocabularies at once: unlimited-per-model; credits per message (Sol 10,
%   Sol Pro 50, deep research 50/task, image 5, voice 5/min); credits per million tokens;
%   rolling windows (Codex Plus 10-100 local messages per 5 hours).
% !! DO NOT quote Anthropic per-plan message counts. They are no longer published and the
%    surviving support article references Sonnet 4 and Opus 4. It is stale.
% ADDED 14 Sept: the Fable ceiling and the buy-more line.
% Max 5x and 20x: Fable is included only within up to 50% of the regular weekly limit.
%   It is a ceiling INSIDE the same allowance, not extra capacity. Past it you can keep
%   using other models, or buy credits to continue on Fable.
% Pro $20: Fable draws pay-as-you-go credits FROM THE FIRST REQUEST, so it is not
%   included at all.
% Two independent clocks: a rolling 5-hour session and a weekly allowance that resets at
%   an account-specific time, visible in Settings > Usage. Waiting out one does not reset
%   the other.
% !! SECONDARY SOURCES (aifreeapi, tokenkarma, techbytes), read 14 Sept. Corroborated
%    across several, not confirmed on Anthropic's own page. Check before stating flatly.
\begin{frame}{Two Different Vocabularies for Usage}
    \begin{itemize}\setlength{\itemsep}{7pt}
        \item Anthropic uses one vocabulary: a rolling five-hour session, plus weekly caps
        \begin{itemize}
            \item Chat, Claude Code and the desktop app all draw from the same pool
        \end{itemize}
        \item OpenAI uses four at once: unlimited per model, credits per message, credits
              per million tokens, and rolling windows
        \begin{itemize}
            \item A single deep research run costs 50 credits; an image costs 5
        \end{itemize}
        \item On the top Anthropic model, half your weekly allowance is the ceiling
        \begin{itemize}
            \item On the \$20 plan it is not included at all: it bills as you go
        \end{itemize}
        \item This is the most confusing part of either product, and the hardest to compare
        \item When you run out you can buy more, by credits or through the API
    \end{itemize}
\end{frame}

% ============================================================
%  PART 2
% ============================================================

% [NEW]

% [NEW] 14 Sept, replacing the multiplier and vetted-tier slides.
% Sources: Ideas/Fall-2026/02_AI_Ecosystems/ecosystems_reference.md section 6, and
% mythos_lsvp.md.
% OpenAI, ChatGPT for Academic Researchers, launched 29 July 2026: 12 months free, $0 at
%   checkout, no auto-renew. Pro-level limits including Sol Pro, expanded deep research,
%   75+ life-science skills. FIVE seats, 1 owner + 4 verified collaborators at the same
%   institution. Eligibility: research faculty OR postdoc at a high-research-activity
%   degree-granting institution, PLUS a paper on arXiv, bioRxiv or ChemRxiv within 3 years.
%   SheerID on institutional email, 1-2 weeks. Excludes API credits, SSO, SCIM.
% OpenAI Researcher Access Program: up to $1,000 API credits, reviewed quarterly
%   (Mar/Jun/Sep/Dec), credits expire in 12 months. NEXT REVIEW IS THIS MONTH.
% Anthropic Team Plan for Scientists: standard seats FREE for 12 months (normally $20),
%   premium seats $15 (normally $100). 1-25 seats. PI-led. Fields: "natural sciences,
%   mathematics, computer science, engineering, and related fields" - "related fields" is
%   UNDEFINED. Capped at 10,000 scientists worldwide. 90 days inactivity reverts pricing.
%   Renews at the then-current price after 12 months. 5-7 business days.
% Anthropic AI for Science: up to $50,000 in credits per project, 6-month term, reviewed
%   the first Monday of every month. API ONLY, not the claude.ai app.
% Anthropic Life Sciences Verification Program: access to the restricted Mythos model for
%   vetted life scientists, US organizations only, first participants enrolled.
% !! NO VERIFIED STUDENT ROUTE. Bruno asked for one; the only sources found were affiliate
%    and SEO sites, not vendor pages. Neither program above admits students directly:
%    OpenAI needs faculty or postdoc status plus a preprint, Anthropic is PI-led and
%    students receive seats through a supervisor's group. If a real student offer exists,
%    verify it on the vendor's own page before adding it.
\begin{frame}{What You Might Not Have to Pay For}
    \begin{center}
    \scriptsize
    \begin{tabular}{p{3.5cm}>{\raggedright\arraybackslash}p{4.2cm}>{\raggedright\arraybackslash}p{4.3cm}}
        \toprule
        Program & Who & What you get \\
        \midrule
        OpenAI, Academic Researchers & Faculty or postdoc, with a preprint in the last three years & Twelve months free at the top limits, five seats \\
        OpenAI, Researcher Access    & Any researcher, reviewed quarterly & Up to \$1,000 in API credits \\
        Anthropic, Team Plan for Scientists & Principal investigators in science and related fields & Free seats for twelve months, \$15 for the larger ones \\
        Anthropic, AI for Science    & Academic and nonprofit projects & Up to \$50,000 in credits per project \\
        Anthropic, Life Sciences     & Vetted life scientists, US institutions & The restricted model, and looser safeguards \\
        \bottomrule
    \end{tabular}
    \end{center}

    \vspace{0.8em}
    \begin{itemize}
        \item None of these is a student route: you get in through a supervisor's group
        \item Both verify against your institutional email, and both take a week or two
    \end{itemize}
\end{frame}

\begin{frame}[plain]
    \vfill
    \begin{center}
        {\Huge Part 2}\\[0.8em]
        {\Large What They Share}
    \end{center}
    \vfill
\end{frame}

% [NEW]
% Parity list from reference section 4. Do not oversell any of these as a differentiator.
\begin{frame}{Both of Them Already Do All of This}
    None of it is a reason to pick one over the other.

    \vspace{0.8em}
    \begin{itemize}\setlength{\itemsep}{5pt}
        \item Projects, memory, deep research, and code execution inside chat
        \item Connectors to your other tools, and MCP support on both sides
        \item Computer use, scheduled tasks, and a million tokens of context
        \item Batch and caching discounts, Office integrations, and a unified desktop app
    \end{itemize}

    \vspace{1em}
    So I will show each of these once, and it will look the same wherever you do it.
\end{frame}

% [STAGE 1]
\begin{frame}{Asking a Question}
    \begin{itemize}
        \item Both will search the web and come back with something about as good as the other
        \item They cite what they used, and you can click through to the source
        \item You can ask follow-ups against what it found, rather than starting again
        \item Everyone already does this, and it is where the two are hardest to tell apart
    \end{itemize}
\end{frame}

% [STAGE 2]
\begin{frame}{Keeping Sources in One Place}
    \begin{itemize}
        \item A Project is a folder with a memory: files, instructions, and its own history
        \item Papers, links and notes go in once, and every later conversation can see them
        \item Standing instructions live there too, so you stop retyping what you care about
        \item If you take one thing away today it is this, and both of them have it
    \end{itemize}
\end{frame}

% [NEW]
% Anthropic: Free 5 projects, paid unlimited. Per-project memory space.
% CAPACITY, added 15 Sept from support.claude.com/en/articles/11473015:
%   RAG is on PAID PLANS ONLY (Pro, Max, Team, Enterprise). It activates automatically
%   "when your project approaches or exceeds the context window limits", with no setup and
%   no user action. Anthropic's claim: "Store up to 10x more content in your projects", and
%   "Response accuracy remains consistent with in-context processing".
%   !! The exact threshold is NOT published; it is tied to the context window for the plan.
%   !! A known bug report (claude-code issue #25759) says files switch to RAG search at 2%
%      capacity because the threshold counts FILES not tokens. Do not quote the 2%: it is a
%      bug report, not documented behavior. Mentioned only so it is not a surprise.
%   !! FREE PLAN GETS NO RAG. A free project simply hits the ceiling.
% OpenAI: unlimited projects on all plans. File caps: Free 5 / Go-Plus 25 / Pro-Business-Edu
%   40. Memory scope is selectable: project-only or default.
%   !! No OpenAI equivalent of the RAG-switch behavior was found. Do not claim symmetry.
\begin{frame}{Projects: The One Real Difference}
    \begin{itemize}
        \item Anthropic caps the number of projects on Free and removes the cap when you pay
        \item OpenAI gives unlimited projects on every plan, and caps the files inside each
        \item Either way a project fills up. Anthropic's then starts searching your files
              instead of reading them all, which holds about ten times as much
        \item So a project is not a bottomless folder. Past a point it retrieves rather
              than reads, and you will not be told when that happened
        \item OpenAI lets you choose whether a project can see your general memory, or only
              its own, which matters if unpublished work lives in one and not another
    \end{itemize}
\end{frame}

% [NEW]
% Batch 50% both. Cache read 90% both. Fast mode 2x price, ~2.5x faster, both.
% Data residency +10% both. Anthropic stacked best case 0.05x base input.
% Cache write differs and is contested: Anthropic 1.25x (5 min) / 2x (1 hr); OpenAI
% reported free, but one source says 1.25x on the GPT-5.6 table. Do not assert.

% [NEW]
% Source: Ideas/Fall-2026/02_AI_Ecosystems/memory.md. Anthropic, 25 Aug 2026.
% Topics added as you chat, not summarized after a conversation ends. Spans Claude chat
%   (web, desktop, mobile) and Cowork, syncing both ways. Stored as editable topic files.
% Settings > Memory: pause, reset, edit or delete individual files, toggle sensitive topics.
% On by default Free/Pro/Max. OFF on Team and Enterprise until an admin enables it.
% Excluded by default, opt-in to include: health, race, ethnicity, religious beliefs,
%   politics, gender identity. Refused outright: government and other sensitive ID numbers,
%   criminal history, immigration status, anything violating the Acceptable Use Policy.
% !! Check OpenAI's current memory scope before implying a difference. Both have memory;
%    the cross-product sync and the editable-file model are the candidate distinctions.
\begin{frame}{What It Remembers About You}
    \begin{itemize}\setlength{\itemsep}{6pt}
        \item Both keep a memory across conversations, not only inside one
        \item Anthropic now writes topics down as you talk rather than summarizing
              afterward, and shows them to you as files you can edit or delete
        \item It spans the chat and the agent, so something you said in one turns up in
              the other
        \item Health, race, religion, politics and gender identity are left out unless you
              opt in
        \item On by default for individuals, and off until an administrator turns it on for
              a team
    \end{itemize}
\end{frame}

% [NEW]
% Anthropic Desktop: Chat / Cowork / Code tabs.
% OpenAI merged the standalone Codex app into the ChatGPT desktop app on 9 July 2026,
% with a Chat / Work / Codex switcher. Six weeks apart. Convergent evolution.
\begin{frame}{Convergent Evolution}
    \begin{itemize}
        \item Both companies built a chat app, a coding agent, and an agent for everything
              that is not code
        \item Both then merged all three into a single desktop application
        \item Anthropic's has Chat, Cowork and Code tabs. OpenAI's has Chat, Work and Codex
        \item They shipped six weeks apart, without copying each other's design
    \end{itemize}
\end{frame}

% ============================================================
%  PART 3
% ============================================================

% [NEW]
\begin{frame}[plain]
    \vfill
    \begin{center}
        {\Huge Part 3}\\[0.8em]
        {\Large Where They Differ in Practice}
    \end{center}
    \vfill
\end{frame}

% [NEW]
% Anthropic Artifacts: interactive React, publishable, 20MB persisted storage.
%   AI-powered artifacts bill to the VIEWER's subscription, not the author's.
% OpenAI Canvas: narrowed to writing and code blocks on 28 May 2026.
\begin{frame}{Artifacts against Canvas}
    \begin{itemize}
        \item Both let the model produce a document you can edit beside the conversation
        \item OpenAI narrowed Canvas to writing and code in May
        \item Anthropic's Artifacts stayed broad: working pages you can publish and share
        \item An artifact that calls the model bills the person viewing it, not the person
              who wrote it
    \end{itemize}
\end{frame}

% [STAGE 3]
% [NEW]
% Source: Ideas/Fall-2026/02_AI_Ecosystems/claude_design.md. Beta as of 6 Aug 2026.
% Chat on the left, canvas on the right. Pro, Max, Team, Enterprise. DISABLED BY DEFAULT
%   on Enterprise. Web and desktop only, no mobile.
% Imports design files, codebases, screenshots. Exports zip, PDF, PPTX, standalone HTML,
%   handoff to Claude Code, plus Adobe, Canva, Gamma, Miro, Replit, Vercel, Wix.
% Draws on the shared usage pool. Known issues: inline comments sometimes do not display,
%   large repos degrade performance, simultaneous multi-person editing is unreliable.
% !! Beta. Check what a Berkeley account actually has before demonstrating this live.
\begin{frame}{A Canvas, Not a Conversation}
    \begin{itemize}\setlength{\itemsep}{6pt}
        \item Anthropic has a separate surface for designing things: chat on one side, a
              canvas on the other
        \item You describe what you want, then refine it by editing the canvas directly
        \item It exports to PowerPoint and PDF, which is the part that matters for this room
        \item It is still in beta, and it is switched off by default on institutional
              accounts
        \item Editing at the same time as somebody else does not reliably work yet
    \end{itemize}
\end{frame}

\begin{frame}{Turning Sources Into Something Readable}
    \begin{itemize}
        \item From a folder of papers, it can write a summary of what changed and why
        \item You can set the shape: length, audience, what to leave out
        \item The same prompt in the two environments gives noticeably different writing
        \item This is the first point where which one you use changes what you get
    \end{itemize}
\end{frame}

% [NEW]
% Anthropic: docx, xlsx, pptx and pdf on ALL plans including Free. 30MB per file.
% OpenAI: Docs/Sheets/Slides native and an Excel add-in, but PowerPoint is ABSENT from the
%   Work desktop flow and PDF is undocumented.
\begin{frame}{Making a File}
    \begin{itemize}
        \item Anthropic writes Word, Excel, PowerPoint and PDF files on every plan,
              including the free one
        \item OpenAI writes to Google Docs and Sheets, and has an Excel add-in
        \item PowerPoint is missing from OpenAI's desktop flow, and PDF output is
              undocumented
        \item If you send colleagues attachments, this matters more than it sounds like it
              should
    \end{itemize}
\end{frame}

% [STAGE 4]  <- cut first if time is short
\begin{frame}{Producing a File Someone Else Can Open}
    \begin{itemize}
        \item It can turn what you have into a document, a slide deck or a spreadsheet
        \item One of them makes the file for you, the other sends you to Google
        \item Where the file lands, and who can open it, differs between them
    \end{itemize}
\end{frame}

% [NEW]
% Claude Code: terminal, VS Code, JetBrains, desktop, web, mobile.
% Codex: web, CLI, IDE, desktop, mobile. Codex is on OpenAI's FREE tier.
\begin{frame}{Two Coding Agents}
    \begin{itemize}
        \item Claude Code and Codex do the same job: they read your files, write, run and
              fix code
        \item Both run in a terminal, in your editor, on the web, and on a phone
        \item Codex is on OpenAI's free tier. Claude Code needs a paid plan
        \item You do not have to write code to use either. Most research use is data
              wrangling and file handling
    \end{itemize}
\end{frame}

% [NEW]
% Cowork: desktop GA, web/mobile beta. ChatGPT Work: GA, and on the free tier.
\begin{frame}{Two Agents for Everything That Is Not Code}
    \begin{itemize}
        \item Cowork and ChatGPT Work are the same idea aimed at documents rather than
              repositories
        \item You describe a job, it works through your files, and you review what came back
        \item ChatGPT Work is generally available and on the free tier
        \item Cowork is generally available on desktop and still in beta on web and mobile
    \end{itemize}
\end{frame}

% [STAGE 5]
\begin{frame}{Sending It Off to Work Alone}
    \begin{itemize}
        \item You can give it a job that takes minutes rather than seconds, and walk away
        \item It decides what to open, what to skip, and when it has enough
        \item It reports back with what it found and where it came from
        \item This is where the two pull apart most
    \end{itemize}
\end{frame}

% [NEW]
% Both support MCP. Anthropic authored it; ~100M monthly downloads; competitors adopted it
% rather than fielding a rival.
\begin{frame}{Connecting Your Own Tools}
    \begin{itemize}
        \item Both can reach your mail, calendar, drive and notes through connectors
        \item Both use MCP, the same open standard, which Anthropic wrote and OpenAI adopted
        \item That is unusual. Competitors normally ship a rival rather than adopt yours
        \item Enabling any connector widens what the model can see considerably
    \end{itemize}
\end{frame}

% [STAGE 6]
% Gmail connector is SCOPED BY LABEL, not by prompt. See session_plan.md.
% Say the enforcement line out loud when doing this: it is the best teaching beat in Part 3.
\begin{frame}{Letting It Read Your Own Material}
    \begin{itemize}
        \item Connectors let it read your mail, calendar, drive and notes
        \item You choose the scope before it runs: one label, one folder, one calendar
        \item Newsletters are the case where this genuinely saves an hour a week
        \item It searches rather than reads everything, so ask what it searched
    \end{itemize}
\end{frame}

% [NEW]
% Anthropic Routines: cron (minimum 1 hour) + authenticated HTTP POST + filtered GitHub
%   webhooks. Research preview. Per-account daily cap.
% OpenAI Tasks: one-off, recurring, and monitoring. HARD COUNT LIMITS: Go 3, Plus 5,
%   Business 10, Pro and Enterprise 15. Max once per hour. Not available in voice, in GPTs,
%   or on Pro models.
\begin{frame}{Work That Runs Without You}
    \begin{itemize}
        \item Both can run a job on a schedule, at most once an hour
        \item OpenAI caps how many you can have: three on Go, five on Plus, fifteen on Pro
        \item Anthropic does not cap the count, and can also trigger on a web request or a
              GitHub event
        \item So one is a scheduler and the other is closer to a small automation system
    \end{itemize}
\end{frame}

% [NEW]
% Source: Ideas/Fall-2026/02_AI_Ecosystems/loops.md. Anthropic, "Getting started with loops",
%   30 June 2026.
% Four kinds: turn-based; goal-based (/goal, stops when the goal is met or max turns);
%   time-based (/loop runs locally, /schedule runs in the cloud); proactive (event or
%   schedule, nobody present).
% Examples given: "/goal get the homepage Lighthouse score to 90 or above, stop after 5
%   tries" and "/loop 5m check my PR, address review comments, and fix failing CI".
% Their line: with a goal loop "Claude doesn't have to make a determination on what is good
%   enough". Guidance: match the interval to how often the thing actually changes.
% Warnings: dynamic workflows can spawn hundreds of agents; a loop that writes code needs a
%   review loop; "a reviewer with fresh context is less biased".
% !! Claude Code only. Neither /goal nor /loop exists in the browser chat. Check OpenAI's
%    equivalent before the parity framing covers scheduled work.
\begin{frame}{Two Kinds of Repeating Job}
    \begin{itemize}\setlength{\itemsep}{6pt}
        \item One kind repeats on a clock: check this every Monday, or every five minutes
        \item The other repeats until a condition is true, with a limit on how many attempts
              it gets
        \item The second is the useful one, because you supply the test instead of the model
              judging its own work
        \item Anthropic's phrasing: this way it does not have to decide what is good enough
        \item Match the interval to how often the thing actually changes, or you pay thirty
              times over for the same answer
    \end{itemize}
\end{frame}

% [STAGE 7]  <- cut first if time is short
\begin{frame}{Having It Run Without You}
    \begin{itemize}
        \item A job that works once can be set to repeat on a schedule
        \item It runs whether or not you are at the machine, and leaves the result waiting
        \item Look at my calendar for tomorrow and tell me what I need to prepare
        \item Check my inbox for anything that actually needs answering today
        \item Watch arXiv for new papers in my area and summarize what is worth reading
    \end{itemize}
\end{frame}

% [NEW]
% Anthropic Skills: open SKILL.md standard (agentskills.io), on EVERY plan including Free,
%   zip-and-upload for non-coders. Plugins + marketplace.json on GitHub/GitLab.
% OpenAI Skills: Business and Enterprise ONLY. Custom GPTs need a paid plan; maintained but
%   strategically stagnant. GPT Store has 1,400+ plugins.
% !! ADDED 15 Sept at Bruno's request: "record a skill". NOT VERIFIED against Anthropic's
%    own documentation. Confirm the feature name and that it is available on the plan you
%    are demoing before saying it out loud.
\begin{frame}{Packaging a Repeatable Job}
    \begin{itemize}
        \item Once a job works, both let you save it so it repeats without rebuilding
        \item Anthropic calls these Skills: a folder with instructions, zipped and uploaded
        \item You can also record one: do the job once while it watches, and it writes the
              skill for you
        \item No code required, and it is an open standard other tools can read
        \item OpenAI's equivalent is Custom GPTs, and its own Skills feature is elsewhere
    \end{itemize}
\end{frame}

% [NEW]
% [NEW]
% Source: Ideas/Fall-2026/02_AI_Ecosystems/output_styles.md, code.claude.com/docs/en/
%   output-styles, plus the comparison table on that page.
% Output styles modify the system prompt (role, tone, format, every turn). CLAUDE.md adds a
%   user message after the system prompt (project conventions, always loaded). Skills load
%   task-specific instructions only when invoked or relevant. --append-system-prompt is a
%   one-off addition. Agents run a subagent with its own system prompt.
% keep-coding-instructions decides whether the built-in software engineering instructions
%   survive; false is right for a writing assistant or data analyst.
% The table below is the CHAT equivalent, which is what this room actually has. Do not
%   present the Claude Code commands.
\begin{frame}{Where an Instruction Should Live}
    \begin{center}
    \small
    \begin{tabular}{lll}
        \toprule
        Put it in & Loaded & Use it for \\
        \midrule
        The prompt           & This turn only        & A one-off request \\
        Custom instructions  & Every turn, always    & Tone and format you always want \\
        Project instructions & Inside that project   & Conventions for one body of work \\
        A saved skill        & When it is relevant   & A procedure you repeat \\
        \bottomrule
    \end{tabular}
    \end{center}

    \vspace{1em}
    \begin{itemize}
        \item The common mistake is keeping everything in the prompt and retyping it forever
        \item The other one is putting everything in the always-loaded layer, which you pay
              for on every single turn
    \end{itemize}
\end{frame}

% This is the asymmetry made concrete: the running demo cannot be finished on both sides.
\begin{frame}{The Asymmetry at the End of the Demo}
    \begin{itemize}
        \item Skills are on every Anthropic plan, including the free one
        \item On OpenAI they are Business and Enterprise only
        \item So the last stage of today's task can be packaged on one side and not the other
        \item And not because of price: there is no personal plan that unlocks it
    \end{itemize}
\end{frame}

% ============================================================
%  PART 4
% ============================================================

% [NEW]
\begin{frame}[plain]
    \vfill
    \begin{center}
        {\Huge Part 4}\\[0.8em]
        {\Large What It Means for You}
    \end{center}
    \vfill
\end{frame}

% [NEW]
% Consumer plans: Anthropic OPT-IN (users chose during the Aug-Oct 2025 rollout).
%   OpenAI OPT-OUT, ON BY DEFAULT. Retention if you accept training: up to 5 years.
% Team/Business, Enterprise/Edu and API: not trained on by default, both sides.
% Private mode: Anthropic Incognito (never trained on regardless of settings);
%   OpenAI Temporary Chats (30-day delete, never trained on).
\begin{frame}{Your Data}
    \begin{itemize}\setlength{\itemsep}{7pt}
        \item On a personal paid or free account, Anthropic asks first. Training is off
              unless you turned it on
        \item OpenAI's is on by default. You have to go and turn it off
        \item If you leave it on, your conversations can be retained for up to five years
        \item On Team, Enterprise, Education and API plans, neither company trains on your
              content
        \item Both have a private mode that is never trained on: Incognito, and Temporary
              Chats
    \end{itemize}
\end{frame}

% [NEW]
% Training choice and retention are SEPARATE controls. Turning training off is not a
% deletion request. Consumer chats stay in the account until the user deletes them;
% deletion then completes from back-end systems within ~30 days, subject to exceptions.
% At OpenAI, training controls, chat history, Temporary Chat and deletion are four
% different things. Do not present "opted out" as "not stored".
\begin{frame}{Turning Training Off Is Not Deleting Anything}
    \begin{itemize}
        \item These are separate controls, and people routinely assume they are one
        \item Opting out of training stops future models learning from your chats
        \item It does not remove the chats. They sit in your account until you delete them
        \item Deletion then works through the back end over about a month, with exceptions
        \item At OpenAI there are four different things here: training, history, temporary
              chats, and deletion
    \end{itemize}
\end{frame}

% [NEW]
% Anthropic trap: a thumbs up/down click sends the ENTIRE conversation into a 5-year store
%   that may be used for training, REGARDLESS of your setting. Also Fable 5 and Mythos 5
%   retain prompts 30 days on every platform, overriding ZDR, no waiver.
% OpenAI trap: memory can retain context beyond what the visible summary shows. Full removal
%   requires deleting across chats, files, connected apps AND the summary. OpenAI states:
%   "Sensitive information may appear in memory if you share it with ChatGPT."
\begin{frame}{Two Traps Worth Knowing}
    \begin{itemize}\setlength{\itemsep}{7pt}
        \item Rating a Claude conversation with thumbs up or down sends the whole
              conversation into a five-year store, whatever your training setting says
        \begin{itemize}
            \item Never rate a conversation that contains unpublished results
        \end{itemize}
        \item ChatGPT's memory can hold more than the summary you are shown
        \begin{itemize}
            \item Clearing it properly means chats, files, connected apps and the summary,
                  separately
        \end{itemize}
    \end{itemize}
\end{frame}

% [NEW]
% [NEW]
% Source: Ideas/Fall-2026/Session_Pool/Teaching_and_Learning/watermarking.md
% Anthropic, 14 Aug 2026, to comply with the EU AI Act; signed the EU Code of Practice on
%   Transparency of AI-Generated Content in July 2026 with ~190 other signatories.
% Method is Google DeepMind's SynthID-Text: the key plus the preceding words settle which of
%   several equally valid word choices is taken. Nothing added, no hidden characters, no
%   extra tokens.
% Returns A PROBABILITY THAT THE TEXT WAS GENERATED BY CLAUDE. It does not tell you whether
%   text is human-written, whether another AI wrote it, or "wrote" vs "heavily edited".
% Limits: short samples fail; sparser on factual passages; barely covers code; a full rewrite
%   removes it; the detection API is in private preview for eligible organizations only.
% NO false positive or false negative rates are published anywhere.
% OpenAI reported to be following; xAI reported not to be; open-weight models cannot be bound.
% !! Verify the OpenAI and xAI positions directly. Full treatment belongs to Session 4.
\begin{frame}{Text Is Starting to Be Watermarked}
    \begin{itemize}\setlength{\itemsep}{6pt}
        \item A watermark is being added to Claude's output, to comply with European law
        \item It works by steering which of several equally good words gets picked, so
              nothing is added and nothing is visible
        \item It answers one question: did Claude probably generate this
        \item It fails on short passages, on factual writing, on code, and on anything
              rewritten
        \item So no watermark found is not evidence a person wrote it, and no error rate has
              been published
    \end{itemize}

    \vspace{0.6em}
    {\small \url{https://www.anthropic.com/news/claude-text-watermark}}
\end{frame}

% Only OpenAI: image generation and editing, full-duplex voice (GPT-Live-1), realtime speech
%   API, transcription and TTS, embeddings API, open-weight models (Apache 2.0), an $8 tier,
%   the GPT Store.
% Only Anthropic: Claude Science (3D protein/genome rendering, HPC over SSH, Modal GPU
%   submission, a reviewer agent that flags bad citations and figures that do not match their
%   code, provenance), Claude Tag in Slack, Claude Design, Agent Teams, Skills on free,
%   Dispatch, Routines' three triggers.
\begin{frame}{What Only One Side Has}
    \begin{itemize}\setlength{\itemsep}{6pt}
        \item Only OpenAI: image generation, live two-way voice, speech and transcription,
              embeddings, and open-weight models you can run yourself
        \item Only Anthropic: Claude Science, which renders structures, submits jobs to a
              cluster, and runs an agent that flags citations and figures that do not match
              the code behind them
        \item Also only Anthropic: Claude Design, an org identity in Slack, and separate
              sessions that can pass findings to each other
        \item If you need images or speech that settles it, and the same goes for job
              submission and provenance
    \end{itemize}
\end{frame}

% [NEW]
% Sora: app ended 26 Apr 2026, API ends 24 Sept 2026, no replacement.
% Fine-tuning: closed to new users, job creation ends 6 Jan 2027, no GPT-5.x support.
% Honest 2026 answer to "can I fine-tune on my corpus": not on a frontier model from either
% vendor. Use prompt caching, Projects, file search, Skills, or self-host gpt-oss.
\begin{frame}{Two Advantages Being Withdrawn}
    \begin{itemize}
        \item Video generation is ending. The app closed in April and the interface closes
              on 24 September, with no replacement
        \item Fine-tuning is closing: no new users now, and it stops entirely in January
        \item So if you were told you could train one of these on your own corpus, that is
              no longer true of either company's frontier models
        \item What is left instead: projects, caching, file search, packaged skills, or
              self-hosting an open model
    \end{itemize}
\end{frame}

% [NEW]
\begin{frame}{Questions to Ask About Your Own Work}
    \begin{itemize}\setlength{\itemsep}{9pt}
        \item Where do you fit in your own workflow? That is what sets the boundary of
              what should be automated
        \item What is actually in your account that you would not want in a training set?
        \item For each tool you would connect, what is the worst thing it could do with
              that access?
        \item What are your guardrails against a confident wrong answer?
        \item How do you check the output without reading every line of it?
    \end{itemize}
\end{frame}

% [NEW]
\begin{frame}{Where to Go From Here}
    \begin{itemize}\setlength{\itemsep}{8pt}
        \item Slides and materials: \url{https://github.com/dlab-berkeley/AI_Pulse}
        \item Workshop suggestions: \url{https://forms.gle/x8KfejgCViE7t6iQA}
        \item What Berkeley already gives you:
              \url{https://life.berkeley.edu/ai-tools-for-berkeley-students}
        \item D-Lab: \url{https://dlab.berkeley.edu}
    \end{itemize}

    \vspace{1em}
    \begin{center}
    Next: AI in the News, 29 September. Then Teaching, Learning and Collaborating,
    13 October. Office hours next week, same day and time, on Zoom.
    \end{center}
\end{frame}

% [NEW]
% Updated 15 Sept with the sources behind the slides added this month. Sources: both vendors'
% pricing and docs pages, the academic program pages, and the dates for Sora and
% fine-tuning shutdowns.
\begin{frame}{References}
    \scriptsize
    \begin{itemize}\setlength{\itemsep}{2.5pt}
        \item Plans, models and prices: \url{claude.com/pricing},
              \url{platform.claude.com/docs}, \url{chatgpt.com/pricing} and OpenAI's API
              pricing page. All checked September 2026
        \item Academic programs: OpenAI, ChatGPT for Academic Researchers, launched
              29 July 2026; Anthropic, AI for Science and Team Plan for Scientists,
              \url{claude.com/programs/team-plan-for-scientists}
        \item Data and training defaults: both vendors' privacy and data-usage pages,
              checked September 2026
        \item Discontinuations: Sora API ends 24 September 2026; fine-tuning job creation
              ends 6 January 2027; Atlas closed 9 August 2026
        \item Berkeley's own tools: \url{life.berkeley.edu/ai-tools-for-berkeley-students}
        \item Auto mode and the 89 percent figure: Anthropic, ``Auto mode is now the default
              in Claude Code'', 7 August 2026. Red teaming by Apollo Research
        \item Permission modes, loops and project retrieval:
              \url{code.claude.com/docs}, \url{support.claude.com}
        \item Memory and watermarking: Anthropic, 25 August and 14 August 2026. The
              watermark uses Google DeepMind's SynthID-Text
    \end{itemize}
\end{frame}

\end{document}
